\scanpage{17}
Caractérisation des foncteurs $(\mathrm{Sch})_{/k}^{o}\longrightarrow\mathrm{Ens}$ qui sont repr

\unsure{Idem|Id.} : \rule{1.7cm}{0.2pt}\quad $\mathrm{Alg}_k\mathrel{-}\mathrm{Ens}$ \rule{2cm}{0.2pt} ].

a) Immersions ouvertes de faisceaux

[stable par chgt de base, par composition]

b) Sup de sous-faisceaux d'un faisceau.

\uline{Prop.} Soit $F$ un foncteur $(\mathrm{Sch})_{/S}\mathrel{-}(\mathrm{Ens})$. Pour que $F$ soit repr. \inserted{(i.e. un \unsure{schéma|schém.})}, il f et s. qu'il satisfasse les conditions suivantes

a) $F$ un faisceau

b) $\exists\ (F_i)_{i\in I}$ famille de \uline{sous faisceaux} \uline{ouverts}, tel que $F=\operatorname{Sup} F_i$ \inserted{(ds cat. de faisceaux)}, et chacun des $F_i$ repr. \unsure{(i.e. un schéma)|par un schéma|i.e. un schéma}

\uline{Cor} On peut \deleted{\unsure{donner|énoncer}} remplacer b) par

b') comme b), les $F_i$ étant de plus des schémas \uline{affines}.

\uline{Cor} (Supposant $S=\operatorname{Spec} k$) Pour qu'un foncteur $k\text{-}\mathrm{Alg}\longrightarrow(\mathrm{Ens})$ soit un faisceau il f et s. que \deleted{\unsure{il soit}} \inserted{ce soit} un \deleted{\unsure{préfais}} \deleted{\unsure{ceau}} faisceau, et qu'on le puisse recouvrir par des sous-faisceaux ouverts affines ....

Ainsi, la notion de schéma est explicitée sans recours à la notion d'espaces topologiques, faisceau d'anneaux, spectre d'un anneau ...

\uline{Applications} a) \uline{Représentabilité du foncteur} $P(E)$

($E$ un module loc. libre sur $S$ -- p. ex. $S=\operatorname{Spec} k$, $E$ un module de \unsure{r.f.|t.f.} sur $k$)

b) \uline{Représentabilité} de \uline{Spec} $A$

c) \uline{Produits fibrés} (donc $\varprojlim$ finies)

% slanted
\margin{sommes sont disjointes et universelles}

d) $\varprojlim$ \inserted{quel.} pour des schémas affines au dessus d'un schéma

% archive sheet number: 16
